% Paper 4, §2 — What we mean by "internal computation": substrate and notation
% Thesis: before any interpretability method can read a model, we must fix
% what is being read. This section codifies the residual-stream substrate,
% the QK/OV decomposition of attention, the FFN as candidate associative
% memory, the LN+W_U readout, and the load-bearing notation
% h^{(l,t)}, d_y, f_i^{(l,t)}, C that the rest of the paper inherits.

\section{What we mean by ``internal computation''}
\label{sec:substrate}

Every interpretability method in this paper reads, intervenes on, or
decomposes a transformer's internal computation. Before the methods can
disagree productively, the substrate they share must be pinned down.
This section codifies four pieces of architectural vocabulary --- the
residual stream as a linear additive bus, attention as a
QK\,$+$\,OV decomposition, the feedforward sublayer as a candidate
associative memory, and the LayerNorm\,$+$\,unembedding as the canonical
readout --- and the load-bearing notation
$\mathbf{h}^{(\ell,t)}, \mathbf{d}_y, f_i^{(\ell,t)}, C$ that
\S\S\ref{sec:logit-lens-family}--\ref{sec:circuits} inherit. Where
Paper~1's architecture chapter establishes the substrate as a design
object, this section recasts it as a \emph{measurement} object: the
quantities the lens reads, the directions the probe projects onto, the
sites the patcher swaps, the subgraphs the circuit picks out.

\subsection{The residual stream as an additive bus}
\label{sec:substrate-residual}

Fix a decoder-only transformer $\mathcal{M}$ with $L$ layers, model
dimension $D$, vocabulary $\mathcal{V}$ of size $V = |\mathcal{V}|$,
$H$ attention heads of head-dimension $d_h$, and unembedding matrix
$W_U \in \mathbb{R}^{D \times V}$ (Paper~1 notation). On an input
sequence $\mathbf{x} \in \mathcal{V}^S$ of length $S$, the model emits
a length-$S$ sequence of residual-stream activations
$\{\mathbf{h}^{(\ell, t)}\}_{\ell=0,\dots,L;\, t=1,\dots,S}$ with
$\mathbf{h}^{(\ell, t)} \in \mathbb{R}^D$. The index $\ell = 0$ denotes
the post-embedding state; $\ell = L$ denotes the final pre-unembed
state. Batched, the residual stream is a tensor of shape $\BSH$, with
$\mathsf{B}$ the batch axis. A head-split per-layer attention tensor is
$\BHSdh$. The output logits are shape $\bsv$.

The pre-LayerNorm transformer block updates the residual stream
additively\footnote{KB excerpt:
\texttt{kb/excerpts/belrose2023-tuned-lens.md\#sec-2-decomposition}.}
\citep[\S 2, Eq.~1]{belrose2023-tuned-lens}:
\begin{equation}
  \mathbf{h}^{(\ell+1, t)}
  \;=\;
  \mathbf{h}^{(\ell, t)} \;+\; F_\ell\!\left(\mathbf{h}^{(\ell, t)}\right),
  \label{eq:residual-update}
\end{equation}
where $F_\ell$ is the residual output of layer $\ell$ --- attention
followed by a feedforward sublayer, each with its own pre-block
LayerNorm or RMSNorm. Eq.~\eqref{eq:residual-update} is the
load-bearing equation of this paper: every component reads
$\mathbf{h}^{(\ell, t)}$, computes some function of it, and writes the
result back into the same residual stream by addition. There is no
gating, no replacement, no nonlinearity at the bus level. Recursing
Eq.~\eqref{eq:residual-update} from any intermediate $\ell$ to the
output gives the residual decomposition\footnote{KB excerpt:
\texttt{kb/excerpts/belrose2023-tuned-lens.md\#sec-2-decomposition}.}
\citep[\S 2, Eq.~2]{belrose2023-tuned-lens}:
\begin{equation}
  \mathcal{M}_{>\ell}\!\left(\mathbf{h}^{(\ell, t)}\right)
  \;=\;
  \mathrm{LN}\!\left[\mathbf{h}^{(\ell, t)}
  \;+\; \underbrace{\sum_{\ell' = \ell}^{L} F_{\ell'}\!\left(\mathbf{h}^{(\ell', t)}\right)}_{\text{residual update}}\right] W_U.
  \label{eq:residual-decomp}
\end{equation}
Here $\mathrm{LN}$ is the model's own final pre-unembed normalization
(LayerNorm or RMSNorm), and $W_U \in \mathbb{R}^{D \times V}$ is the
unembedding (often weight-tied with the input embedding). Two
identifications follow immediately. Setting the residual update to
zero in Eq.~\eqref{eq:residual-decomp} gives the \emph{logit lens}
of \S\ref{sec:logit-lens-family} \citep[\S 2,
Eq.~3]{belrose2023-tuned-lens}; replacing
$\mathbf{h}^{(\ell, t)}$ with $A_\ell\mathbf{h}^{(\ell, t)} +
\mathbf{b}_\ell$ before the same readout gives the \emph{tuned
lens}. Both are read-offs of the same bus.

\begin{intuition}
The residual stream is a linear additive bus. Attention heads, MLPs,
and LayerNorms are devices wired into that bus, each reading a
projection of $\mathbf{h}^{(\ell, t)}$ and writing back a vector that
\emph{adds} to it; nothing replaces or multiplies the running state.
This is what makes lenses cheap (the bus's contents at any layer are
already in the basis the unembedding wants, up to a per-layer affine
correction); what makes patching causally clean (substituting a single
addend leaves all other addends untouched in expectation); and what
makes the dictionary view of \S\ref{sec:sparse-autoencoders}
applicable (a sum of sparse contributions decomposes into a sparse
linear combination of feature directions in $\mathbb{R}^D$). Returning
to canonical form: Eq.~\eqref{eq:residual-update} says every layer is
$\mathbf{h}_{\text{new}} = \mathbf{h}_{\text{old}} + \Delta$, and every
interpretability method is a particular way of reading or perturbing
the running sum.
\end{intuition}

\subsection{Attention as QK\,$+$\,OV}
\label{sec:substrate-attention}

A single attention head $h$ at layer $\ell$ contributes to $F_\ell$ a
term that the \emph{Mathematical Framework for Transformer
Circuits}~\deepencite{elhage2021-mathematical-framework HTML at
transformer-circuits.pub, citation-pending} factors into two
read--write circuits. The \textbf{QK circuit} computes the
attention-weight pattern from the residual stream:
\begin{equation}
  \alpha^{(\ell, h)}_{t, s}
  \;=\;
  \softmax_{s'}\!\left(\frac{1}{\sqrt{d_h}}\,
    \big(W_Q^{(\ell, h)}\,\mathbf{h}^{(\ell, t)}\big)^\top
    \big(W_K^{(\ell, h)}\,\mathbf{h}^{(\ell, s')}\big)
  \right)\!\bigg|_{s' = s},
  \label{eq:qk}
\end{equation}
with $W_Q^{(\ell, h)}, W_K^{(\ell, h)} \in \mathbb{R}^{d_h \times D}$.
The \textbf{OV circuit} determines what is read and written along
each attended position:
\begin{equation}
  \mathrm{Attn}_{\ell, h}\!\left(\mathbf{h}^{(\ell, t)}\right)
  \;=\;
  W_O^{(\ell, h)}\sum_{s = 1}^{S}
    \alpha^{(\ell, h)}_{t, s}\,
    W_V^{(\ell, h)}\,\mathbf{h}^{(\ell, s)},
  \label{eq:ov}
\end{equation}
with $W_V^{(\ell, h)} \in \mathbb{R}^{d_h \times D}$ and
$W_O^{(\ell, h)} \in \mathbb{R}^{D \times d_h}$. The QK pair determines
\emph{where} a head looks; the OV pair determines \emph{what} it
copies and \emph{which direction} that copy writes into the residual
stream. Multi-head attention is a sum over heads of
Eq.~\eqref{eq:ov}. The decomposition is load-bearing for circuit
analysis: when \citet{wang2022-ioi} report (\S\ref{sec:circuits})
that a Name Mover Head ``copies the attended name'', the
human-interpretable claim is exactly that its OV circuit acts as a
near-identity on the residual-stream subspace encoding ``name'', and
its QK circuit attends to the IO position; the
attention-pattern--vs--effect split is the QK--vs--OV split.

\subsection{Feedforward sublayers as candidate associative memory}
\label{sec:substrate-ffn}

The other half of $F_\ell$ is the feedforward (or MLP) sublayer:
\begin{equation}
  \mathrm{MLP}_\ell\!\left(\mathbf{h}^{(\ell, t)}\right)
  \;=\;
  W_{\text{down}}^{(\ell)}\,\sigma\!\left(W_{\text{up}}^{(\ell)}\,\mathbf{h}^{(\ell, t)}\right),
  \label{eq:mlp}
\end{equation}
with $W_{\text{up}}^{(\ell)} \in \mathbb{R}^{d_{\text{ff}} \times D}$,
$W_{\text{down}}^{(\ell)} \in \mathbb{R}^{D \times d_{\text{ff}}}$, an
elementwise nonlinearity $\sigma$ (GELU / SiLU / SwiGLU per the model;
biases and gating dropped here for brevity), and intermediate width
$d_{\text{ff}} \approx 4D$ in dense models. The hypothesis that
ties this to interpretability is the \emph{Localized Factual
Association} reading of \citet{meng2022-rome}: the rows of
$W_{\text{up}}^{(\ell)}$ act as \emph{keys} and the columns of
$W_{\text{down}}^{(\ell)}$ act as \emph{values} in an associative
memory indexed by subject representations, with mid-layer MLPs
storing the bulk of factual recall content. This is the framing the
ROME rank-one MLP edit of \S\ref{sec:rome} acts on; its empirical
warrant comes from the activation-patching localization at
$\mathrm{AIE}_{\mathrm{MLP}} = 6.6\%$ vs.\
$\mathrm{AIE}_{\mathrm{attn}} = 1.6\%$ on GPT-2 XL
\citep[\S 2.2]{meng2022-rome}\footnote{KB excerpt:
\texttt{kb/excerpts/meng2022-rome.md\#sec-2-2}.}. We adopt the
key--value framing as the \emph{anchor for §10}, not as a settled
fact: alternative framings (FFNs as feature-direction \emph{detectors}
in superposition, the \citet{leask2025-sae-not-canonical} non-atomic
critique of \S\ref{sec:sae-canonical}) coexist in the 2026 literature
and are taken up where they bear on the methods.

\subsection{LayerNorm and the canonical readout}
\label{sec:substrate-readout}

Two normalization variants appear in the residual-stream literature
covered here. LayerNorm~\citep{ba2016-layernorm} centers and rescales
$\mathbf{h}^{(\ell, t)}$ along its $D$-dim axis with a learned affine;
RMSNorm drops the centering. Either is denoted $\mathrm{LN}[\,\cdot\,]$
when the distinction does not matter for the argument, and pre-block
or post-block placement is folded into the model's specification.
The canonical \emph{readout} is the model's own pre-unembed
$\mathrm{LN}$ followed by $W_U$; the layer-$L$ logits are
\begin{equation}
  \mathrm{logits}(\mathbf{x})_t \;=\; \mathrm{LN}\!\left[\mathbf{h}^{(L, t)}\right] W_U \;\in\;\mathbb{R}^V.
  \label{eq:readout}
\end{equation}
Equation~\eqref{eq:residual-decomp} is exactly Eq.~\eqref{eq:readout}
re-expressed at an arbitrary depth $\ell$. The lens family of
\S\ref{sec:logit-lens-family} reuses Eq.~\eqref{eq:readout} verbatim
on $\mathbf{h}^{(\ell, t)}$ for $\ell < L$; the
correlation-vs-causation gap of \S\ref{sec:correlation-causation}
is precisely the gap between ``Eq.~\eqref{eq:readout} \emph{would}
produce token $r$ at layer $\ell$'' and ``the model \emph{uses} the
$\ell$-th-layer information that drives that prediction''.

\subsection{What we mean by a feature}
\label{sec:substrate-feature}

Throughout this paper a \textbf{feature} is a direction
$\mathbf{d} \in \mathbb{R}^D$ in residual-stream space (or in some
other activation space named explicitly), \emph{not} a neuron and
\emph{not} necessarily a basis vector. The motivating commitment is
the \emph{linear representation hypothesis}~\deepencite{elhage2021-mathematical-framework
§residual-stream linear-representation, citation-pending}: many
human-interpretable properties of the input are linearly readable
from the residual stream, so the natural object of interpretability
is an axis (or a low-dimensional subspace) of $\mathbb{R}^D$ rather
than an individual coordinate. Every method in this paper picks
features under a different commitment.
\begin{itemize}
  \item \S\ref{sec:logit-lens-family} (lens). The \emph{token-aligned}
    features are the columns of $W_U$ themselves: the lens projects
    $\mathbf{h}^{(\ell, t)}$ onto every column simultaneously and
    reads the $\arg\max$.
  \item \S\ref{sec:probing} (probe). A \textbf{property direction}
    $\mathbf{d}_y \in \mathbb{R}^D$ is a feature associated with an
    external label $y$ --- truth, sentiment, syntactic role,
    sycophancy --- discovered either by supervised fitting
    (Eq.~\eqref{eq:linear-probe}-equivalent) or in closed form by the
    mass-mean rule
    $\mathbf{d}_{\mathrm{MM}} = \boldsymbol{\mu}_+ - \boldsymbol{\mu}_-$.
  \item \S\ref{sec:sparse-autoencoders} (SAE). A \textbf{dictionary
    feature} $i$ at layer $\ell$, position $t$ is the activation
    $f_i^{(\ell, t)} \in \mathbb{R}_{\geq 0}$ of the $i$-th latent in
    a sparse autoencoder trained on $\mathbf{h}^{(\ell, t)}$, with
    associated decoder direction $\mathbf{d}_i$. The dictionary
    $\{\mathbf{d}_i\}_{i=1}^M$ is overcomplete ($M \gg D$); on any
    given input, $\mathbf{f}^{(\ell, t)} \in \mathbb{R}^M$ is sparse.
\end{itemize}
A feature is therefore an \emph{axis}, not a \emph{coordinate}: the
neuron basis is one possible feature basis, and the production
record~\deepencite{elhage2021-mathematical-framework §superposition,
citation-pending} suggests it is rarely the right one. The choice of
basis --- $W_U$'s columns, a learned $\mathbf{d}_y$, an SAE's
$\mathbf{d}_i$, or some other --- is the methodological commitment
that distinguishes one interpretability technique from the next.

\subsection{What we mean by a circuit}
\label{sec:substrate-circuit}

A \textbf{circuit} $C \subseteq \mathcal{M}$ is a subgraph of
$\mathcal{M}$'s computational graph --- nodes are forward-pass
components (residual-stream sites, attention heads, MLP sublayers,
SAE/transcoder features) and edges are the residual / attention /
projection links between them --- whose function $C(\mathbf{x})$
agrees with $M(\mathbf{x})$ on the task distribution after \emph{knockout}
of $\mathcal{M} \setminus C$~\citep[\S 2]{wang2022-ioi}\footnote{KB
excerpt: \texttt{kb/excerpts/wang2022-ioi.md\#sec-2-definition}. Wang
et al.\ define knockout via mean ablation on a reference distribution
$p_{\mathrm{ABC}}$ rather than zero ablation, on the empirical
grounds that ``zero ablation [\ldots] lead[s] to noisy results''.}.
A circuit \emph{claim} commits to three quantitative criteria from
\citet[\S 4]{wang2022-ioi}\footnote{KB excerpt:
\texttt{kb/excerpts/wang2022-ioi.md\#sec-4-criteria}.}:
\begin{description}
  \item[Faithfulness.] $C$ performs the task as well as the whole
    model: $C(\mathbf{x}) \approx \mathcal{M}(\mathbf{x})$ on the task
    distribution under a chosen metric.
  \item[Completeness.] $C$ contains every node used to perform the
    task: knocking out $\mathcal{M} \setminus C$ does not change
    behavior beyond noise.
  \item[Minimality.] $C$ contains no task-irrelevant nodes: each
    node in $C$ fails an ``ablate-it-and-see-if-it-matters'' test.
\end{description}
The substantive content of a circuit claim is therefore threefold:
the relevant components are \emph{localized} (a small fraction of the
graph), the explanatory edges are \emph{causal} (knockout matters),
and the subgraph is \emph{sparse} (most of $\mathcal{M}$ is
irrelevant on the task distribution). \citet{wang2022-ioi} note that
real circuits routinely pass faithfulness and minimality but fail the
strongest completeness tests --- the structural fact that
\S\ref{sec:circuit-ioi} unpacks via the \emph{backup head}
phenomenon, in which ablating Name Mover heads reactivates latent
backups that were dormant in the unablated forward pass. The
features-in-circuits dichotomy is fundamental: a circuit can route a
feature from one position to another without itself \emph{computing}
that feature, so the wiring diagram (the circuit) and the encoding
(the feature) are independent specifications of an
explanation~\citep[\S 2.1]{wang2022-ioi}.

\subsection{The notation table}
\label{sec:substrate-notation}

Table~\ref{tab:substrate-notation} fixes the symbols inherited by the
rest of the paper. Tensor shapes are listed in the
$\mathsf{B} \!\times\! \mathsf{S} \!\times\! \mathsf{D}$ convention of
Paper~1.

\begin{table}[h]
\centering
\small
\begin{tabularx}{\linewidth}{l l X}
\toprule
\textbf{Symbol} & \textbf{Shape / type} & \textbf{Meaning} \\
\midrule
$\mathbf{h}^{(\ell, t)}$ & $\mathbb{R}^D$ ($\BSH$ batched) & residual-stream activation at layer $\ell \in \{0,\dots,L\}$, position $t \in \{1,\dots,S\}$ \\
$F_\ell$ & $\mathbb{R}^D \to \mathbb{R}^D$ & residual contribution of layer $\ell$: $\mathrm{Attn}_\ell + \mathrm{MLP}_\ell$ (Eq.~\eqref{eq:residual-update}) \\
$\mathrm{LN}[\,\cdot\,]$ & $\mathbb{R}^D \to \mathbb{R}^D$ & LayerNorm or RMSNorm, model-specific \\
$W_U$ & $\mathbb{R}^{D \times V}$ & unembedding matrix; canonical readout is $\mathrm{LN}[\mathbf{h}]\,W_U$ \\
$\alpha^{(\ell, h)}_{t, s}$ & scalar in $[0,1]$ & attention weight from query $t$ to key $s$, head $h$, layer $\ell$ (Eq.~\eqref{eq:qk}) \\
$W_Q^{(\ell, h)}, W_K^{(\ell, h)}, W_V^{(\ell, h)}$ & $\mathbb{R}^{d_h \times D}$ & QK and OV input projections \\
$W_O^{(\ell, h)}$ & $\mathbb{R}^{D \times d_h}$ & OV output projection (writes back into the bus) \\
$\mathbf{d}_y$ & $\mathbb{R}^D$ (unit-norm) & property direction associated with label $y$ (probe / steering target) \\
$f_i^{(\ell, t)}$ & $\mathbb{R}_{\geq 0}$ & SAE / transcoder feature $i$'s activation at $(\ell, t)$ \\
$\mathbf{d}_i$ & $\mathbb{R}^D$ & decoder column for feature $i$ (the feature's direction in residual space) \\
$C$ & subgraph of $\mathcal{M}$ & circuit, in the Wang-2022 sense \\
$B, S, D, H, d_h, V, L$ & --- & batch, sequence, model-dim, heads, head-dim, vocab, layers (Paper~1) \\
\bottomrule
\end{tabularx}
\caption{The notation inherited by §3--§11. $\mathbf{h}^{(\ell, t)}$ is
the residual-stream activation site every method reads;
$\mathbf{d}_y$ is what supervised probes find; $f_i^{(\ell, t)}$ is
what unsupervised SAEs find; $C$ is what circuit reverse-engineering
returns. The tensor-shape conventions $\BSH$, $\BHSdh$, $\bsv$
follow Paper~1's architecture chapter.}
\label{tab:substrate-notation}
\end{table}

The four objects in the second half of the table fix what each method
\emph{produces}. A lens (\S\ref{sec:logit-lens-family}) maps
$\mathbf{h}^{(\ell, t)}$ to a layer-$\ell$ token distribution via
Eq.~\eqref{eq:residual-decomp}. A probe (\S\ref{sec:probing}) reads
$\mathbf{h}^{(\ell, t)} \mapsto \sigma(\mathbf{d}_y^\top
\mathbf{h}^{(\ell, t)})$ for a learned property direction
$\mathbf{d}_y$. An SAE (\S\ref{sec:sparse-autoencoders}) decomposes
$\mathbf{h}^{(\ell, t)}$ into a sparse $\mathbf{f}^{(\ell, t)} \in
\mathbb{R}^M$ with components $f_i^{(\ell, t)}$ along directions
$\mathbf{d}_i$. Activation patching (\S\ref{sec:patching}) intervenes
on $\mathbf{h}^{(\ell, t)}$ (or on a finer-grained component) and
measures the indirect effect on a downstream metric. Circuit
reverse-engineering (\S\ref{sec:circuits}) returns a subgraph $C$
satisfying the Wang criteria. All four methods read or write the same
underlying object.

\subsection{Forward link: four reads of one bus}
\label{sec:substrate-forward}

The unifying picture of this section is that the residual stream
defined by Eq.~\eqref{eq:residual-update} is one object, and the five
interpretability methods this paper covers are five different ways of
reading or perturbing it.
\S\ref{sec:logit-lens-family} reads the bus's running prediction by
applying the canonical readout~\eqref{eq:readout} at every layer.
\S\ref{sec:probing} reads the bus under supervision by training a
classifier that projects $\mathbf{h}^{(\ell, t)}$ along a learned
$\mathbf{d}_y$. \S\ref{sec:sparse-autoencoders} reads the bus
\emph{without} supervision by fitting a sparse dictionary
$\{\mathbf{d}_i\}_{i=1}^M$ such that
$\mathbf{h}^{(\ell, t)} \approx \sum_i f_i^{(\ell, t)} \mathbf{d}_i$.
\S\ref{sec:patching} \emph{intervenes} on the bus, swapping a
component's contribution to $F_\ell$ between a clean and a corrupted
forward pass and measuring the downstream metric change.
\S\ref{sec:circuits} returns a subgraph of $\mathcal{M}$ whose nodes
are precisely the bus sites and the components writing into them.
The methods commit to different objects --- the $W_U$ columns, the
property direction $\mathbf{d}_y$, the SAE direction $\mathbf{d}_i$,
the component $C$ at a knockout site, the subgraph $C$ as a wiring
diagram --- but the underlying read site is the same
$\mathbf{h}^{(\ell, t)}$ throughout. The convergences and divergences
of \S\ref{sec:cross-method} live exactly at the boundaries between
these reads of one bus.
